\section{Discussion}

The main value of the EEG companion study is temporal triangulation. ARC-ADHD already defines a behavioral event sequence and a monitoring-anchored recovery endpoint. EEG can test whether a second, independently measured signal changes coherently around the same controlled event without becoming part of the core intervention.

The most important design choice is to avoid interpreting conventional frequency bands as direct mental-state labels. ADHD EEG research shows task-related oscillatory differences at the group level, but heterogeneity remains substantial and current evidence does not support simple diagnostic translation \parencite{michelini2022eventoscillations,slater2022eegsubtypes}. Accordingly, the present design emphasizes within-child event locking, target-rule-matched references, reproducibility, artifact sensitivity, and incremental value.

Muse S Athena may be operationally attractive because it combines wearable EEG with a form factor intended for repeated use \parencite{interaxonAthena2026}. That convenience does not remove the need to verify raw Athena research access, montage details, timing, and signal-quality behavior. A technically convenient stream that cannot be synchronized or audited would not support the proposed event-level inference.

A positive result would justify a later study asking whether EEG-informed adaptation or neurofeedback improves decisions beyond behavioral telemetry. It would not itself establish therapeutic benefit. A negative result would be equally useful because it would argue against adding a burdensome sensor to the core application.
